\documentclass[11pt,a4paper]{article}
\usepackage[a4paper,margin=20mm]{geometry}
\usepackage{amsmath,booktabs,longtable,graphicx,xcolor}
\usepackage{xurl}
\usepackage[colorlinks=true,urlcolor=blue,linkcolor=blue]{hyperref}
\usepackage{fancyhdr}
\setlength{\emergencystretch}{4em}
\setlength{\parskip}{0.4em}
\setlength{\parindent}{0pt}
\renewcommand{\arraystretch}{1.15}
\pagestyle{fancy}
\fancyhf{}
\fancyhead[L]{AI/Human Music Detection: Optimization Supplement}
\fancyfoot[C]{\thepage}
\setlength{\headheight}{15pt}
\title{Interpretable AI/Human Music Detection\\\large External Transfer and Source-Held-Out Optimization}
\author{Thesis experiment archive}
\date{7 September 2026}
\begin{document}
\maketitle
\begin{abstract}
This English thesis supplement separates accepted external-transfer findings
from a subsequent, prospectively frozen development experiment. The completed
transfer comparison shows that additional signal descriptors can improve one
class endpoint while worsening the other, and that group weighting can reverse
an apparent improvement. The expanded development experiment contains 2,207
eligible excerpts and evaluates all 127 combinations of seven feature families
at five training caps. Its real performance is not yet numerically accepted at
this checkpoint. Synthetic implementation tests are not music-detection
benchmarks. The previously released thesis PDF is preserved unchanged.
\end{abstract}
\textbf{Reading boundary.} Section 1 reports accepted measurements. Section 2
documents the frozen experiment, not a new accuracy claim. Section 3 reports
input coverage and interpretation risks. Paths in the supplement are relative
to \path{artifacts/audio_phenomena_expansion_20260907/} unless stated otherwise.

\textbf{Feature key.} S: separated-vocal high-frequency spectrum;
D: non-vocal dynamics; R: detected rhythm variation;
P: section-duration/bar-count variety; F: phase and group delay;
H: pitch-class trajectories; M: long-range, transposition-aware recurrence.
These are descriptive proxies, not direct measurements of authorship or intent.

\section{Matched external AI and Human transfer}
\label{sec:saraga-mureka-frozen-transfer}

The same 3,175 previously saved development models were applied without
adaptation to 500 source-labelled Mureka recordings and 103 source-labelled
Saraga Hindustani Human recordings. The grid contains all 127 nonempty subsets
of S/D/R/P/F/H/M, five training caps and five saved folds. Model identities,
candidate identifiers, cap/fold assignments and training/test identity hashes
match exactly across sources. All original transformations and the raw-score
threshold 0.5 remain fixed. No fitting, threshold tuning or feature selection
uses either external cohort.

Table~\ref{tab:saraga-mureka-allcap} reports five-model means at the all cap,
which uses all available training groups within each original fold, not one
fit on the entire 1,604-row development cohort. Actual training sizes range
from 1,250 to 1,328 rows. Repeated model applications to the same recordings
are dependent; model ranges are descriptive, not confidence intervals.

\begin{table}[htbp]
\centering
\small
\begin{tabular}{lrrrr}
\toprule
Families & AI recall & Human spec. & Human FPR & Group-equal spec. \\
\midrule
S+D+R+P & 76.76 & 67.77 & 32.23 & 75.18 \\
S+D+R+P+F & 63.76 & 85.05 & 14.95 & 87.06 \\
S+D+R+P+H & 80.08 & 67.57 & 32.43 & 76.84 \\
S+D+R+P+M & 83.28 & 59.61 & 40.39 & 65.56 \\
All seven & 79.36 & 80.58 & 19.42 & 76.45 \\
\bottomrule
\end{tabular}
\caption{Source-separated frozen transfer, all cap. Rates are percentages.
AI recall concerns Mureka only; Human specificity/FPR concern Saraga only.
Group-equal specificity averages five performer/release component rates.}
\label{tab:saraga-mureka-allcap}
\end{table}

Adding F changes Mureka sensitivity by $-13.00$ percentage points and Saraga
specificity by $+17.28$ points. Adding M changes these endpoints by $+6.52$ and
$-8.16$ points. Both trade-off directions persist across all five cap means.
Adding H changes the all-cap endpoints by $+3.32$ and $-0.19$ points; its
directions depend on cap. All seven improve the all-cap means by $+2.60$ and
$+12.82$ points, but do not improve both endpoints at every cap. These are
matched feature-set model differences, not causal effects of an isolated
audio intervention.

Single-family all-cap Mureka sensitivity/Saraga specificity are S
89.24/33.01, D 48.32/67.57, R 91.00/64.66, P 67.12/69.71, F 64.20/90.29,
H 80.56/71.84 and M 82.12/17.28 percent. In particular, M alone falsely labels
82.72 percent of these Human recordings as AI on average. The results caution
against universal interpretation of recurrence or high-frequency signatures;
source, genre and production effects are not isolated by this design.

\paragraph{Component sensitivity.}
Saraga has five connected performer/release components of sizes 72, 14, 10,
5 and 2. All-seven versus S/D/R/P all-cap specificity changes are respectively
$+16.39$, $+10.00$, $+8.00$, $-8.00$ and $-20.00$ percentage points. Hence the
recording-weighted $+12.82$-point gain becomes only $+1.28$ points with equal
component weights. The smallest groups have substantial sampling uncertainty.
At caps 25, 50, 100 and 200, recording-weighted Human specificity improves
while equal-component specificity declines. Neither aggregation should be
silently substituted for the other, and five components are not five datasets.

\paragraph{Coverage and limits.}
All 103 Human recordings remain present. S/D/F/H/M are complete on 103,
R on 98 and P on 77. Missing values use only stored training transformations.
Coverage mechanisms are documented in
\path{SARAGA103_MEASUREMENT_COVERAGE_EN.md}. Provenance labels do not provide
forensic certainty; source and class are confounded in these external tests.
Saraga is outside the frozen training cohort but has historical research
exposure through earlier 10-second inputs. Exact file-hash nonoverlap does
not prove work, performance or pretraining independence. No pooled accuracy,
ROC AUC, calibrated probability, significance or universal winner is claimed.

\paragraph{Reproducibility and next development boundary.}
Independent numerical auditing verified all 327,025 Saraga scores, with maximum
absolute error $1.7208456881689926\times10^{-15}$ and exact agreement of every
threshold decision. The audit covers all model/component integer counts and
635 aggregate cells. Parent checks additionally reconciled 63,190 presentation
and source numerical values. All 11 Saraga result files (177,911,960 bytes)
are hash-verified locally and on NFS.

Code is \path{code/score_saraga103_frozen_v4.py},
\path{code/audit_saraga103_transfer_v1.py} and
\path{code/present_saraga_mureka_frozen_transfer_v1.py}. The score freeze is
\path{preregistration/saraga103_frozen_v4_transfer_frozen_v1.json}.
Predictions and per-model results are in
\path{results/saraga103_frozen_v4_transfer_results_v1/}; all 635 matched cells,
60 fixed overview rows, 100 paired differences and component tables are in
\path{results/presentation_saraga_mureka_frozen_transfer_v1/}. Acceptance is
\path{audit/saraga_mureka_transfer_parent_acceptance_v1.json}; the numerical
audit is \path{audit/saraga103_frozen_v4_transfer_numerical_audit_v1.json}.
The narrative is \path{SARAGA_MUREKA_FROZEN_TRANSFER_RESULTS_EN.md}.

These external scores are now exposed. A later optimized development version
must preserve this experiment and explicitly label any inclusion of these
603 recordings as development reuse. Combined with the original cohort this
would yield 2,207 observations, 1,311 Human and 896 AI; it would not create an
untouched confirmation set. Broader source/generator-heldout testing is needed
before claiming improved generalization.

\clearpage
% Requires booktabs. Methodology supplement; not a completed-results claim.
\section{Exploratory source-held-out development after external evaluation}
\label{sec:v5-exploratory-sourceheldout}

Following the frozen Mureka and Saraga transfer experiments, a new, explicitly
exploratory development package combines 2,207 accepted 60-second excerpts:
1,311 Human and 896 AI recordings, from six Human sources and two generators.
This consists of the original 1,604 development recordings, 500 Mureka
recordings and 103 Saraga recordings. It is the eligible common-context subset,
not the entire collected corpus. The original 438 locked and 50 pilot
recordings remain excluded. A separate origin ledger preserves original
roles, metadata and admission flags; no upstream role is overwritten.

Mureka and Saraga outcomes have already been examined. Their reuse cannot be
presented as untouched confirmation of hypotheses motivated by those outcomes.
Source-held-out evaluation prevents the specified within-fit source overlap,
but does not remove research-level exposure or establish pretraining
independence. The preceding frozen experiments remain unchanged.

\subsection{Complete family and quantity ablation}

All 54 existing descriptors are retained: S (15), D (3), R (3), P (6),
F (15), H (6), and M (6). Every recording and its missing values remain in
the common cohort. The primary design evaluates all 127 nonempty subsets at
five caps, with training-only median imputation and missingness indicators.
Two additional all-cap diagnostic modes use median-imputed values without
indicators, or missingness indicators alone. These diagnostics do not select
features or models.

With seed 20260907, global groups are assigned to five hash buckets. Each
source-held-out arm excludes the entire held-out source and every current
test-bucket group from training. Its test set contains current-bucket examples
from the held source and the opposite class. The verified schedule contains
28 Human-source folds, 10 generator folds, and five ordinary group-disjoint
descriptive folds. Saraga occupies only three buckets, with 77, 14 and 12
Human recordings; its two empty source-fold cells are recorded, not filled.

A cap limits groups per class/source within the training fold, retaining all
recordings in selected groups. It is not a recording count or a balanced
total sample size. Deterministic group rankings are nested across caps.
Actual training-row ranges over the 43 valid folds are shown in
Table~\ref{tab:v5-caps}.

\begin{table}[htbp]
\centering
\caption{Verified training quantities in the prospective v5 schedule.
The cap counts groups per training source, not songs.}
\label{tab:v5-caps}
\begin{tabular}{lrr}
\toprule
Group cap & Minimum training rows & Maximum training rows \\
\midrule
25 & 194 & 389 \\
50 & 355 & 678 \\
100 & 578 & 898 \\
200 & 963 & 1,297 \\
All & 1,351 & 1,805 \\
\bottomrule
\end{tabular}
\end{table}

The fixed weighted ridge model has penalty 10 and an unpenalized intercept.
Training weights balance classes, sources within class, groups within source,
and recordings within group, then sum to training size $n$. Consequently,
the mean-loss regularization is $10/n$: quantity also changes effective
regularization. This design preserves the historical model implementation,
but is not a causal quantity-only or diversity-only experiment. All decisions
use raw, unbounded identity-link score $\geq 0.5$; these scores are not
calibrated probabilities. No full-cohort refit or threshold search is performed.

\subsection{Aggregation and acceptance}

Within each held-source arm, candidate, cap, mode and source, the recording-level
class-conditional correct rate is the correct count divided by unique
out-of-fold recordings: Human specificity or AI sensitivity, not joint accuracy.
The component-level endpoint is the unweighted mean of individual global-group
correct rates. Averaging fold-level component means would be incorrect when
buckets contain different numbers of components, as in Saraga's 2/1/2 split.

Two-class summaries average source-pair metrics within a fold, valid folds
within a held-source arm, and then held-source arms within a fold type.
Human-source, generator and ordinary group holdouts remain separate.
Pooled-fold companion metrics are explicitly distinguished. No pooled
out-of-fold AUC is computed across separately fitted score scales, and no
raw predictions are pooled across unlike source-held-out arms. All 635
primary subset/cap cells per reporting arm are retained; no winner or
statistical significance is inferred from descriptive fold ranges.

The frozen schedule specifies 27,305 primary and 10,922 diagnostic fits,
with 10,687,558 predictions in total. Every training transformation, model,
test identity and output metric is persisted. At the time this methodology
supplement was written, execution was in progress: neither completion nor
independent numerical acceptance is claimed. Final performance reporting
requires a complete committed publication and independent score and metric
reconstruction.

\subsection{Reproducibility references}

All paths below are relative to
\path{artifacts/audio_phenomena_expansion_20260907/}:
\begin{itemize}
\item Protocol: \path{EQUAL60_EXPLORATORY_V5_PROTOCOL_EN.md}.
\item Derived-data code: \path{code/prepare_evaluation_inputs_v5.py};
evaluation code: \path{code/evaluate_new_phenomena_v5.py}.
\item Inputs and original-role ledger:
\path{results/equal60_exploratory_package_v5_v1/}.
\item Parent input acceptance:
\path{audit/exploratory_v5_package_parent_acceptance_v1.json}.
\item Reviewed split schedule:
\path{results/equal60_exploratory_v5_draft_v1/}.
\item Execution authorization:
\path{preregistration/equal60_exploratory_v5_frozen_v1.json}.
\item In-progress outputs, not yet accepted results:
\path{results/equal60_exploratory_v5_results_v1/}.
\item Existing third-generator context constraints:
\path{THIRD_GENERATOR_CONTEXT_ELIGIBILITY_EN.md}.
\end{itemize}

\clearpage
\section{Input coverage and dependence diagnostics}
\label{sec:v5-coverage-balance}

Numerical completeness is assessed per family: every descriptor must be finite.
Blank and NaN entries are missing, not zero. The input package contains exactly
matching metadata and feature identities in the same order.

\begin{table}[htbp]
\centering\small
\begin{tabular}{lrrrrrrrr}
\toprule
Source & Rows & S & D & R & P & F & H & M \\
\midrule
MTG-Jamendo & 481 & 481 & 481 & 475 & 388 & 479 & 478 & 471 \\
Mureka v9 & 500 & 500 & 500 & 500 & 485 & 500 & 500 & 499 \\
Suno & 396 & 396 & 396 & 395 & 354 & 396 & 396 & 395 \\
MAESTRO v3 & 300 & 300 & 300 & 298 & 194 & 300 & 300 & 300 \\
MedleyDB & 156 & 156 & 156 & 156 & 128 & 155 & 156 & 156 \\
MoisesDB & 238 & 238 & 238 & 238 & 223 & 238 & 237 & 234 \\
Saraga Hindustani & 103 & 103 & 103 & 98 & 77 & 103 & 103 & 103 \\
URMP & 33 & 33 & 33 & 33 & 17 & 33 & 33 & 33 \\
\bottomrule
\end{tabular}
\caption{Complete-family counts in the expanded 60-second package. Counts do
not certify that an estimator measures its intended musical construct.}
\end{table}

P completeness is 97.0\% for Mureka, 64.7\% for MAESTRO and 51.5\% for URMP.
A classifier may therefore exploit estimator availability instead of, or in
addition to, section-length values. The frozen all-cap missingness-only and
values-only diagnostic arms address this concern without dropping incomplete
recordings. Partial-family counts are retained in
\path{EXPLORATORY_V5_COVERAGE_AND_GROUP_BALANCE_EN.md}.

S remains numerically complete even on instrumental datasets. Finite features
from a separated vocal estimate do not establish the presence of a singer:
instrument residuals can enter the estimate. A future vocal-presence or
stem-reliability gate would require external validation and a new prospective
comparison; it is not introduced into the running frozen experiment.

The 1,779 recorded global groups are unevenly distributed: MTG 481, Mureka 500,
Suno 396, MAESTRO 49, MedleyDB 78, MoisesDB 237, Saraga 5 and URMP 33.
Singleton provenance keys do not prove independent artists or works. Saraga's
largest component contains 72 of 103 recordings. Thus collecting more files
from an already dominant component need not improve diversity, and both
recording-weighted and equal-component endpoints remain necessary.

\paragraph{Context constraints for a third generator.}
The stored HeartMuLa and ACE-Step cohorts contain no eligible native 60-second
examples. Short excerpts must not be padded to manufacture long-range
recurrence evidence. The available long DiffRhythm examples remain protected
pilot data. A separate short-context branch would also have to preserve shared
conditioning groups: HeartMuLa/ACE share 500 prompt groups, while 481 Udio
groups overlap current MTG Human groups. The inventory and eligibility limits
are recorded in \path{THIRD_GENERATOR_CONTEXT_ELIGIBILITY_EN.md}. No additional
generation or acquisition is reported here.

\paragraph{Implementation acceptance, not detection performance.}
Twenty-four auditor/presenter unit tests passed on the NFS CPU runtime.
A separate full-grid synthetic publication passed independent checking of
17,780 model instances and 56,896 predictions, including transforms, residuals,
decisions and metrics. Maximum score error was
$5.773159728050814\times10^{-15}$; decisions matched exactly.
This is numerical interoperability evidence, not real-music performance.

Evidence: \path{audit/equal60_v5_synthetic_numerical_audit_v1.json}.
Code: \path{code/audit_equal60_v5_results.py} and
\path{code/present_equal60_v5_results.py}.
The synthetic directory \path{results/v5_synthetic_interop_l4ayhwmw/}
must never enter thesis music-detection performance tables.
\end{document}
